% 第 6 章 总结与展望
%===============================================================================
% 修改记录:
%   2026-05-05  创建
%===============================================================================

\chapter{总结与展望}
\label{ch:conclusion}

\section{工作总结}
\label{sec:summary}

本文围绕 STM32F407 平台的 FSMC 外设展开研究，主要工作与结论如下：

\begin{enumerate}
    \item 系统分析了 FSMC 的内部结构与 Bank 地址映射，明确了 SRAM 异步模式下
          \texttt{ADDSET}、\texttt{DATAST}、\texttt{ADDHLD} 等关键时序参数与
          系统时钟的约束关系，得到访问周期的解析表达式（式~\ref{eq:twr}）。
    \item 提出了一种自适应时序参数搜索方法，从厂家推荐值出发逐步收紧关键
          参数，并以写读校验为停机判据，无需人工查表即可获得稳定可用的
          最小周期配置。
    \item 搭建了基于 STM32F407 + IS62WV51216 SRAM + ILI9341 LCD 的实验
          平台，验证表明本文方法可将 SRAM 写带宽从 35.2~MB/s 提升至
          58.7~MB/s，LCD 全屏刷新延迟从 28~ms 降至 18~ms。
\end{enumerate}

\section{未来展望}
\label{sec:future}

本文工作仍可在以下方向延伸：

\begin{itemize}
    \item \textbf{DMA 加速}：本文实测均基于 CPU 直接读写。引入 DMA 后可释
          放 CPU，进一步提升整机吞吐，需重新分析 AHB 矩阵的仲裁影响。
    \item \textbf{多 Bank 并行}：现有方法仅优化单一 Bank。在 SRAM + LCD
          同时工作时，两 Bank 之间的访问顺序与时序裕量需联合优化。
    \item \textbf{温度自适应}：将温度传感器读数作为搜索算法输入，在不同
          工况下动态调整时序，提升老化器件的兼容性。
\end{itemize}
